% STATUS: ACTIVE -- companion to formC_ugap.tex. Everything that is not an
%          equation: what the coordinate buys, what it does not buy, the
%          measurements, and the open items. Independently verified 2026-08-13;
%          three errors in the main file were found and fixed before any code
%          was written against it (see "Verification" below).
% ======================================================
% Compile: xelatex formC_ugap_ref.tex
% Path:    reference/eq17_analysis/derivation/formC_ugap_ref.tex
% ======================================================

\documentclass[11pt,a4paper,fontset=none]{ctexart}
\usepackage[fleqn]{amsmath}
\usepackage{amssymb,bm}
\setlength{\mathindent}{0pt}
\usepackage{geometry}
\geometry{margin=2.0cm}
\setlength{\parindent}{0pt}
\setlength{\parskip}{0.80em}
\linespread{1.08}
\newcommand{\ba}{\bar{a}}
\newcommand{\bw}{\bar{w}}
\begin{document}

\section*{Why this coordinate}

Any autonomous scalar flow $\mathrm{d}\ba/\mathrm{d}\bw = g(\ba)$ is
straightened by $u = \int\!\mathrm{d}\ba/g(\ba)$. Here $g = (1-\ba)^{2}(1+\delta a)$
gives $u = 1/(1-\ba)$ and $\mathrm{d}u/\mathrm{d}\bw = 1+\delta a$. The height
writing's law is NOT autonomous ($\ba'$ depends explicitly on $\bw$), so it has
no such coordinate; the state writing's defect and its remedy are two sides of
the same property.

Physically $u = (1+\delta a)(\bw-\bw_0)$ is the gap to the implied wall. The
state changes from ``how strong is the gain'' to ``how far am I'', and distance
propagates by exactly how far one moved.

\subsection*{What it buys}

\begin{enumerate}
\item The one-signed quadrature ratchet is removed IDENTICALLY, not reduced.
$\mathrm{d}u/\mathrm{d}\bw$ is constant in $\bw$, so the one-step integral is
exact. Measured on the canonical command trajectory (same $\Sigma$, $\ba$-Euler
versus exact-$u$): $+1.641\,\%$ at the end of the descent, $+2.122\,\%$ at the
end of the oscillation, strictly one-signed throughout, and exactly first order
in $\Delta t$ ($+1.077 / +0.542 / +0.272\,\%$ at $\Delta t/2, /4, /8$). Closed
form: the $\ba$-Euler step equals $u[k+1] = u[k] + \Sigma + \Sigma^{2}/u +
\mathcal{O}(\Sigma^{3})$, and $\Sigma^{2}/u > 0$ regardless of direction --- that
is the ratchet. $+2.1\,\%$ is the right size against the $2.8$ point gap between
arm (a) at $4.68\,\%$ and the true-slope arm at $1.87\,\%$.
\item The predict increment reads $\hat{u}$ nowhere: $\mathrm{d}u =
(1+\delta a)\mathrm{d}\bw$. That is the property the ``$\ba'$ from the TRUE
height'' arm had, obtained here without knowing the true height.
\item $u$ is affine in the unknown wall position, so a Gaussian posterior is
closer to the true posterior shape in $u$ than in $\ba$. An EKF in $u$ is a
DIFFERENT, non-equivalent estimator --- this is the mechanism by which it can
win.
\item $(u,\delta a)$ separate cleanly: a wall-position error is a constant
offset in $u$, a $\delta a$ error is a ramp proportional to accumulated
displacement ($F_e(4,5) = \hat{\Sigma}$).
\end{enumerate}

\subsection*{What it does not buy}

No information. $u = (1+\delta a)(\bw-\bw_0)$ with $\bw$ commanded and free, so
$(u,\delta a) \leftrightarrow (\bw_0,\delta a)$ is a known bijection at fixed
$\bw$, and Fisher information is invariant under smooth reparametrisation. What
changes is conditioning. (That invariance assumes the prior is pushed forward
consistently --- see the seed section.)

The self-amplification is NOT abolished. The term $-2(1-\hat{\ba})\Sigma$ is
exactly the coordinate's own time-varying Jacobian: a fixed $e_u$ really does
correspond to an $e_{\ba}$ that grows $123\times$ between $\bar h = 22.2$ and
$2.0$. What changes is that it moves from the propagation Jacobian to the
output map $\partial\ba/\partial u = u^{-2}$, where it does not compound.

The structural limits of the $\ba$ writing survive unchanged: in a hold
$\hat{\Sigma} = 0$ so $F_e(4,5) = 0$ and $\delta a$ is unobservable through the
state path; over a closed oscillation $\sum\hat{\Sigma} = 0$ so the state path
carries zero net information about it. Only $y_2$ sees $\delta a$ continuously.

Also unchanged: the $y_1$ gain-correction path through $K_1 = PH_1^{\top}/S_1$
(measured to supply essentially all of the across-seed spread, and to be
correctly priced at $0.04$--$16\,\%$ of the realised optimum); the MA(2)
augmentation the controller runs but no derivation describes ($\approx 21\,\%$
of the spread); $\bar h_{\mathrm{safe}} = 1.5$ as an underived house value.

\section*{Seed, prior, and a well-posedness constraint}

$\delta a$ enters the LAW here, so its prior enters $u$ at $k=0$ and
$P_{u\,\delta a}[0] \neq 0$. With $\hat u[0] = 22.22$ and $\sqrt{P_{\delta a}} =
0.5$ the three contributions to $P_{uu}[0]$ are $123.5$ ($\delta a$), $0.012$
(wall position) and $\hat u[0]^{4}\mathrm{floor}_a^{2}$; the disturbance term
dominates the wall term by $10^{4}$.

$\mathrm{floor}_a$ MUST be re-derived for THIS family. The single-curve family
of \texttt{formC\_state\_dist} declares $13.05\,\%$ RMS / $15.34\,\%$ sup; this
family, which carries $\delta a$, is the one declared in
\texttt{formC\_state\_da} at $0.66\,\%$ / $4.65\,\%$. The choice matters far
more here than in $\ba$ coordinates because the floor is amplified by
$\hat u[0]^{4}$:

\begin{equation*}
\begin{array}{lccc}
\mathrm{floor}_a & \sigma_u & P(u<1) & \text{status} \\[0.2em]
0.0306\ (\text{envelope sup}) & 15.11 & 8.0\,\% & \text{ill-posed} \\
0.0056\ (\text{local at seed}) & 2.77 & \approx 0 & \text{well-posed}
\end{array}
\end{equation*}

$u>1 \Leftrightarrow \ba<1$, and $u \to 1$ is the singularity of the output map
$1-1/u$. Eight percent of prior mass at $\ba<0$ sitting on that singularity is
not a width question. In $\ba$ coordinates the same two choices differ only by a
factor $5.5$ in width; a linear push-forward cannot preserve both tails. The
implementation therefore uses the local floor and clamps $u \geq u_{\min}$
(with $\bar h_{\mathrm{safe}} = 1.5$ the natural value is $u_{\min} = 1.5$),
and clamps $\hat v$ likewise.

Reporting $\hat{\ba} = 1-1/\hat u$ carries a Jensen bias $\approx
-P_{uu}/\hat u^{3}$: negligible once converged ($0.01\,\%$ at $u=2$,
$0.03\,\%$ at $u=4.2$) but $2.18\,\%$ at the seed under the sup floor versus
$0.07\,\%$ under the local floor --- the same choice again.

\section*{Verification, 2026-08-13}

Independent check before implementation. Central claim, S6, S7 and S8 confirmed
term by term, including a time-varying similarity transform
$F^{u}[k] = T[k+1]F^{\ba}[k]T[k]^{-1}$ with $T = \hat u^{2}$ that reproduces the
printed $F_e$ exactly. Three errors were found and are fixed in the main file:

\begin{enumerate}
\item $P_{uu}[0]$ omitted the $\delta a$ term and $P_{u\,\delta a}[0]$ was set to
zero. Inherited from the wrong sibling (see 3), where the zero is correct
because $\delta a$ does not enter the law there. Magnitude-dominant.
\item The delay back-off was labelled exact; it drops the
$(\delta\bw_3-\delta\bw_1)$ channel, which \texttt{formB\_ws\_ref} already
declares as a truncation. Measured contribution $0.5$--$3\,\%$ of $\sqrt{R_2}$
(peak $11\,\%$), i.e.\ $\leq 0.1\,\%$ of variance, but it is correlated with the
states, so it is bias and not variance. Now declared. Note $R_2$'s $d\,\hat
v^{-4}Q_{uu}$ exists precisely because $H$ omits that channel, so ``exact'' and
that term cannot both stand.
\item The sibling cross-reference pointed at \texttt{formC\_state\_dist.tex},
whose $\delta a$ is a different object (additive per step, absent from the law).
The $\ba$-coordinate version of THIS law is \texttt{formC\_state\_da.tex}.
\end{enumerate}

Also reported: \texttt{formC\_state\_da.tex} is internally inconsistent on
$Q_{33}$ (prefactor $4k_BT/(a_oR)$ in one place, $4k_BT a_o/R$ in its $\xi$);
the main file here is the correct one.

\section*{Pre-registration}

On the canonical scenario over the 8 house seeds, overall relative gain-error
RMS: the same estimator in $\ba$ coordinates gives $4.68 \pm 2.56\,\%$; fed
$\ba'$ from the true height it gives $1.87 \pm 0.93\,\%$; production tier-2
gives $2.07 \pm 1.38\,\%$.

\begin{equation*}
\begin{array}{ll}
\text{lands near } 4.7 & \text{the gap is information; the coordinate cannot help} \\
\text{lands near } 1.9 & \text{the gap is conditioning} \\
\text{lands below } 1.9 & \text{also expected: the true-slope arm still carries the}\\
 & \text{Euler ratchet, which this writing removes as well}
\end{array}
\end{equation*}

Arm (a) against $4.68\,\%$ is the clean controlled comparison: same law, same
estimator, same measurement chain, differing only in the quadrature and in
whether the integrand reads the state.


\section*{Result, 2026-08-13}

Eight house seeds, paired, canonical scenario. Overall relative gain-error RMS:

\begin{equation*}
\begin{array}{lccc}
 & \text{desc peak} & \text{osc RMS} & \text{overall RMS} \\[0.2em]
\ba \text{ coordinate, arm (a)}      & 15.25\pm6.30 & 4.24\pm1.74 & 4.59\pm2.56 \\
u \text{ coordinate, arm (a)}        & 16.02\pm7.70 & 4.42\pm2.07 & 5.46\pm2.90 \\
\ba \text{ coordinate, arm (b)}      & 17.52\pm6.41 & 7.39\pm7.32 & 7.33\pm4.61 \\
u \text{ coordinate, arm (b)}        & 17.09\pm4.80 & 4.22\pm2.09 & 4.76\pm2.71 \\
\ba' \text{ from the true height}    & 3.23\pm0.98  & 1.58\pm1.01 & 1.87\pm0.93 \\
\text{production tier-2}             & 4.22\pm1.80  & 1.90\pm1.00 & 2.07\pm1.38
\end{array}
\end{equation*}

Arm (a) paired difference $u-\ba$: $+0.874 \pm 1.605$, $t = +1.54$, $u$ better on
$3/8$. THE FIRST PRE-REGISTERED BRANCH: it lands near $4.7$, not near $1.9$. The
coordinate does not close the gap, and the point estimate is on the wrong side
of zero. The $2.4\times$ gap to production tier-2 is INFORMATION, not
conditioning.

That is decisive because the ratchet removal is not in doubt --- it was measured
at $+2.1\,\%$, one-signed, and this writing removes it identically. Removing it
buys no accuracy. So the ratchet was already being corrected by the
measurements, and what survives is what the true-slope arm removes and this
coordinate cannot: the integrand reads the state, and no smooth
reparametrisation changes what the data say about the state. The invariance
argument of ``What it does not buy'' is thereby confirmed by experiment rather
than only asserted.

Read off the figures: the error trace has the SAME SHAPE in both coordinates ---
same features at the same times --- shifted bodily. Seed 7 shifts the wrong way
(final hold $-3 \to -5\,\%$), seed 101 the right way (final hold returns toward
zero). A shift whose sign is set by the seed, not by the coordinate, is what
$t = +1.54$ over $3/8$ means.

Two findings that are not about the headline:

\begin{enumerate}
\item In arm (b) the coordinate HELPS: $7.33 \to 4.76$, $u$ better on $7/8$
($t = -1.71$). And $\delta\hat{a}$ actually moves here --- $-0.145 \pm 0.058$,
one-signed $8/8$ --- where in \texttt{formC\_state\_dist} it never left noise
($-5.4\times10^{-8}$). The difference is structural: $\delta a$ enters the law
in this family, so $F_e(4,5) = \hat\Sigma$ makes it observable through
displacement. The disturbance state is only estimable in the writing where it
means something. It still does not reach arm (a), so this rehabilitates the
mechanism without rehabilitating the arm.
\item The well-posedness table was right, and it was worth having. The envelope
floor gives $15.32\,\%$ overall RMS against the local floor's $5.46\,\%$, and
clamps $u$ on up to $23.5\,\%$ of steps against $4.0\,\%$. Predicted before the
run from $P(u<1) = 8.0\,\%$ alone.
\end{enumerate}

The clamp is not the explanation for the headline: per-seed clamp count against
the paired difference correlates at $-0.008$.

\section*{Standing after this result}

Three independent routes have now failed to close the same gap --- an additive
disturbance state, a multiplicative one, a uniform-in-time covariance inflation,
and a coordinate change that provably removes the discretisation error. The one
thing that does close it is knowing where the wall is. Production tier-2 does
not know either, and matches the true-slope arm, so the remaining question is
not how to integrate the gain but whether anything in the state writing can
substitute for evaluating $\ba'$ at a height. On the present evidence the answer
is no, and the state writing's case has to rest on something other than accuracy
on this scenario --- extrapolation beyond the traversed band being the obvious
candidate, and still unmeasured.

\end{document}
